\section{Setup}
\label{sec:setup}

\textbf{Task.} 100 problems drawn from a decontaminated KodCode subset (n-gram containment
against EvalPlus). \textbf{All 50 full-length arms in this paper use the identical 100 items}
(verified: unique item-set count = 1), so every cross-arm comparison is paired.

\textbf{Protocols.}
\emph{ReAct} --- a text protocol with \texttt{Thought:} / \texttt{Action:\ run\_tests} / \texttt{Action\ Input:\ \textasciigrave{}\textasciigrave{}\textasciigrave{}python\ \ldots{}\textasciigrave{}\textasciigrave{}\textasciigrave{}}
/ \texttt{Observation:}; the tool name is fixed in the template.
\emph{Function calling (FC)} --- OpenAI-style, \texttt{tool\_choice:\ auto}, one tool \texttt{run\_tests} taking a
single \texttt{code} string. Schema variants: \textbf{terse} (bare \texttt{\{"type":"string"\}}), \textbf{rich}
(parameter description, counter-example, code sample), \textbf{strict} (\texttt{strict:\ true},
\texttt{additionalProperties:\ false}).

\textbf{Serving.} vLLM 0.27.1~\citep{kwon2023vllm,vllm_toolcalling_docs}, per-family parser and chat template as documented.
\texttt{temperature=0} is greedy but \textbf{not bit-wise deterministic} under continuous batching:
identical prompts can differ across batch compositions. The main table is a single sample
per item; \S\ref{sec:variance} quantifies sampling variability separately at \texttt{temperature=0.6} rather than
assuming greedy decoding removes it.
\texttt{max\_tokens\ =\ 2048} for \textbf{both} protocols (see \texttt{ERRATA.md} \S\ref{sec:intro} --- an earlier asymmetry
favoured FC; re-running under the true shared limit changed final pass by $\leq$1 item).

\textbf{Measured quantities.} For each item's first turn we record: whether the server parsed a
tool call; the raw output; the tool name called; the raw \texttt{arguments}; \texttt{finish\_reason}; and
a protocol-agnostic intent classification (\S\ref{sec:intent}). Execution, per-turn results, and final
pass are recorded per turn.

\subsection{Intent criterion}
\label{sec:intent}

\emph{The criterion's source, the three tiers, and every
prompt and tool schema used in this paper are reproduced verbatim in Appendix~\ref{app:prompts}.}

Reporting only server-parsed calls is the practice under examination, so we need an
independent measure of what the model emitted. A single classifier (\texttt{analysis/intent.py})
is shared by all text, tables and figures:

\begin{itemize}
\tightlist
\item
  \textbf{tight} --- a JSON object naming \texttt{run\_tests} with an \texttt{arguments} object whose \texttt{code} is a
  \textbf{string literal containing real Python}. Excludes two contaminants we observed:

  \begin{enumerate}
  \def\labelenumi{(\roman{enumi})}
  \tightlist
  \item
    echoing the injected schema (\texttt{parameters}/\texttt{description} present, no \texttt{arguments});
  \item
    illustrative pseudo-code (\texttt{\{"code":\ test\_code\}} --- an identifier, not a string).
  \end{enumerate}
\item
  \textbf{strong} --- trajectory tag \texttt{json\_named\_call} or \texttt{xml\_tool\_call}.
\item
  \textbf{weak} --- JSON-shaped fragments only; at 1.5B these are almost entirely false positives
  and we report them separately rather than merging them.
\end{itemize}

\subsection{Human validation of the intent criterion}
\label{sec:humanval}

The scale trend in \S\ref{sec:scale} rests entirely on an automated classifier, so we validated it
against human judgement. 98 outputs were sampled \textbf{stratified by classifier verdict}
(40 \texttt{tight}, 28 \texttt{strong}-but-not-\texttt{tight}, 30 neither) --- deliberately over-sampling the
decision boundary so both false positives and false negatives are estimable. Sampling is
seeded and the annotator saw only the raw output, never the classifier's label.

We report \textbf{two rounds}, because the first one is itself a finding.

{\small
\begin{longtable}[]{@{}lllllll@{}}
\caption{Human validation of the intent criterion. Two annotation rounds on a stratified sample; round 2 supplied the model's own preceding turn as context. The classifier is used only to count what the server discarded, so its recall matters more than its precision --- and after correction it misses almost nothing.}\label{tab:validation}\\
\toprule
& agreement & Cohen's $\kappa$ & precision & recall & FP & FN \\
\midrule
\endfirsthead
\toprule
& agreement & Cohen's $\kappa$ & precision & recall & FP & FN \\
\midrule
\endhead
Round 1 (raw output only) & 86.7\% & 0.713 & 70.0\% & 96.6\% & 12 & 1 \\
After adjudication & \textbf{96.9\%} & \textbf{0.936} & \textbf{95.0\%} & 97.4\% & 2 & 1 \\
\bottomrule
\end{longtable}
}

The 12 round-1 false positives split into two causes, and the second one is ours, not the
annotator's. \textbf{Eleven are reading-order failures}: the tool call sits mid-document,
after a fenced \verb|```python| block (match positions 1209--2585 in outputs of median length
\textasciitilde2000); the annotator read the code block, concluded ``this is a direct answer,
not a call'', and stopped. \textbf{One is an instrument failure}: for a 32B item the matched
span begins at character 2585 and the annotation pack truncated every output at 2400, so the
evidence was \emph{not in the pack at all}. That item could not have been labelled correctly.
Re-shown the matched span alone, 10 of 13 disputed items were reversed.

We flag the second cause rather than folding it into the first because \textbf{the annotation
instrument reproduced the very failure this paper is about}: a truncating tool silently
removed the evidence, and what came out the other side was ``the model did not call the
tool''. The corrected pack (below) carries every output in full.

We report this rather than only the adjudicated figure because \textbf{it reproduces the paper's
mechanism on a human reader}: the evidence was present in every case, and what determined
whether it was seen was the order in which the text was scanned. The serving parser is
blocked by format; the annotator was blocked by reading order. Both produce the same
conclusion --- ``the model did not call the tool'' --- from the same bytes.

\textbf{The first adjudication was one-sided, and we do not treat its $\kappa$ as a reliability
estimate.} Auditing the procedure after the fact: the 13 adjudicated items were selected
\emph{exactly} as the items where the annotator disagreed with the classifier (13/13
overlap), and all 10 reversals moved \emph{toward} the classifier, none away. Under that
selection rule agreement can only rise, so $0.713 \to 0.936$ is a property of which items
were re-examined, not evidence of reliability. We report the second annotator instead.

\subsubsection*{A second independent annotator, and what it revealed about the first}

A second annotator, blind to the classifier, to the first annotator's labels and to the
first pack's item order, labelled the identical 98 items under a verbatim-identical rubric.
Two changes were made to the instrument: outputs are carried \textbf{in full} (the first
pack truncated at 2400 characters), and the adjudication rule was fixed in advance to review
\emph{all} three-way disagreements rather than only those disagreeing with the classifier.

\begin{table}[htbp]
\centering
\caption{Inter-annotator reliability. Restricting to items where both annotators read
identical bytes is not a convenience: 11 of the 12 raw disagreements are items the first
pack truncated, where A1 answered ``no call'' because the call had been cut from their copy.}
\label{tab:interannotator}
\small
\begin{tabular}{@{}lcc@{}}
\toprule
Cohen's $\kappa$ (95\% CI, bootstrap) & all items ($n$=97) & \textbf{identical bytes ($n$=62)} \\
\midrule
A1 vs.\ A2 (inter-annotator) & 0.736 [0.590, 0.863] & \textbf{0.967 [0.897, 1.000]} \\
A2 vs.\ classifier           & 0.936 [0.855, 1.000] & \textbf{0.967 [0.897, 1.000]} \\
A1 vs.\ classifier           & 0.712 [0.559, 0.847] & 0.934 [0.834, 1.000] \\
\bottomrule
\end{tabular}
\end{table}

\textbf{The apparent unreliability was the instrument, not the raters.} Of the 12 items where
the two annotators disagree, \textbf{11 are items whose full text exceeds 2400 characters},
and in every one of them A1 answered \texttt{N} while A2 answered \texttt{Y} --- A1 said ``no
call'' because the call was not in the copy they were given. On the 62 items where both read
identical bytes, the two annotators agree on 61 and \textbf{inter-annotator $\kappa$ = 0.967}.

We keep this rather than quietly regenerating a clean number, because it is the paper's own
thesis arriving uninvited in our methods section: \textbf{a truncating instrument sat between
the evidence and the observer, and what came out was ``the model did not call the tool''} ---
the same sentence the \texttt{hermes} parser produces, from the same bytes, for the same
reason. We were measuring censoring with a censored instrument.

\textbf{Third-party adjudication, and the final labels.} Fifteen items were not unanimous
across A1, A2 and the classifier. All fifteen went to a third adjudicator under the
pre-registered rule --- selected by \emph{any} three-way disagreement, never by disagreement
with the classifier --- who saw full text and no one else's labels. The adjudicator sided
with A2 on 11, with the classifier on 9, and with A1 on 6, and \textbf{moved away from the
classifier on 6 of 15}. That last number is the check that matters: the first round's
adjudication moved away from the classifier zero times out of thirteen, which is what an
outcome looks like when the selection rule has already decided it.

Combining the 83 unanimous items with the 15 adjudicated ones gives a gold standard that at
no point consulted the classifier.

\begin{table}[htbp]
\centering
\caption{Validation against the adjudicated gold standard. The inter-annotator column is
restricted to items where both annotators read identical bytes, because A1's pack truncated
at 2400 characters; the classifier and the gold standard are unaffected by that truncation
and are scored on the full sample.}
\label{tab:validation2}
\small
\begin{tabular}{@{}lcc@{}}
\toprule
Cohen's $\kappa$ (95\% CI, bootstrap) & full sample & identical bytes \\
 & ($n$=97) & ($n$=62) \\
\midrule
A1 vs.\ A2 (inter-annotator) & 0.736 [0.595, 0.866] & \textbf{0.967 [0.893, 1.000]} \\
Classifier vs.\ gold standard & \textbf{0.871 [0.756, 0.958]} & 0.934 [0.829, 1.000] \\
\bottomrule
\end{tabular}
\end{table}

\textbf{Correction factor.} Against the gold standard the \texttt{tight} stratum's precision
is 36/40 = \textbf{0.900}, so \S\ref{sec:scale}'s headline 80/100 at 32B corrects to \textbf{$\approx$72}.
The criterion errs in both directions --- 4 over-counts against 2 calls it misses entirely,
for a net over-count of 2 --- so the correction is smaller than the precision figure alone
suggests. \textbf{The trend is unaffected}: applying 0.900 uniformly gives
$0, 3.6, 18.9, 32.4, 72$ against a server-side zero at every size. We report uncorrected
classifier counts in all tables with this factor stated, rather than rescaling silently.

For the record, this number moved twice as the validation improved, and we report the
trajectory rather than only its endpoint: 0.950 under the first, one-sided adjudication;
0.975 under A2's labels alone; \textbf{0.900} against the adjudicated gold standard. The
first was inflated by its selection rule, the second rested on a single rater, and only the
third is built from a rule fixed before the labels were seen.

\textbf{A bound on the measurement itself.} The probe stores at most 4000 characters of raw
output per item; 10 of the 98 sampled outputs reach that ceiling, and one of the grey-zone
items above is a call truncated by it. The classifier and both annotators read identical
bytes, which is what agreement requires, but a call emitted beyond character 4000 is
invisible to all three. Emitted-call counts are accordingly a \textbf{lower bound} ---
conservative for our claim, since it can only understate the censoring we report.

\subsection{Validity gating}
\label{sec:gating}

An arm is admissible for pass-rate comparison only if it has exactly 100 items, exit code 0,
\textbf{zero request errors}, and provenance (model / protocol / temperature / \texttt{max\_tokens} /
adapter) matching the intended configuration. Arms with request errors retain a valid
\textbf{error-rate census} but are marked N/A for pass rates and \emph{no p-value is computed for
them}. This rule is enforced in code (\texttt{validate\_arms.py}, \texttt{final\_table.py}), not by
convention: an earlier draft computed \emph{p}=0.648 on such an arm, and that number is
retracted.

\begin{center}\rule{0.5\linewidth}{0.5pt}\end{center}
